% !TeX root = manual_fundamentos_matematicas.tex
% ------------------------------------------------------------
% Capítulo 7
% ------------------------------------------------------------

\chapter{Los números enteros}

\begin{objetivos}
  \item Introducir el conjunto de los números enteros.
  \item Estudiar sus operaciones y su orden.
  \item Presentar el valor absoluto.
  \item Introducir la división euclídea.
\end{objetivos}

\section{Motivación: la necesidad de ampliar los naturales}

En el capítulo anterior se introdujo el conjunto de los números naturales, que denotamos por
\[
  \mathbb{N}=\{0,1,2,3,\ldots\}
\]
o bien por
\[
  \mathbb{N}^{*}=\{1,2,3,\ldots\},
\]
según se quiera incluir o no el número cero. En este manual adoptaremos la convención
\[
  \mathbb{N}=\{0,1,2,3,\ldots\}.
\]

Los números naturales permiten contar objetos y ordenar cantidades discretas. Sin embargo, presentan una limitación esencial: no siempre es posible resolver ecuaciones de la forma
\[
  a+x=b,
\]
donde \(a\) y \(b\) son números naturales y \(x\) es la incógnita.

Por ejemplo, si queremos resolver
\[
  5+x=3,
\]
no existe ningún número natural \(x\in\mathbb{N}\) que satisfaga la igualdad. Para poder resolver este tipo de ecuaciones es necesario ampliar el sistema numérico e introducir números que representen diferencias negativas.

\begin{ejemplo}
La ecuación
\[
  7+x=10
\]
tiene solución en \(\mathbb{N}\), concretamente \(x=3\). En cambio, la ecuación
\[
  10+x=7
\]
no tiene solución en \(\mathbb{N}\). Para resolverla se introduce el número \(-3\), de manera que
\[
  10+(-3)=7.
\]
\end{ejemplo}

\begin{advertencia}
El símbolo \(-3\) no debe entenderse simplemente como una resta pendiente de realizar. En el conjunto de los números enteros, \(-3\) es un número con entidad propia: el opuesto de \(3\).
\end{advertencia}

\section{El conjunto de los números enteros}

\begin{definicion}[Números enteros]
El conjunto de los números enteros, denotado por \(\mathbb{Z}\), está formado por los números naturales, sus opuestos y el cero:
\[
  \mathbb{Z}=\{\ldots,-3,-2,-1,0,1,2,3,\ldots\}.
\]
\end{definicion}

La letra \(\mathbb{Z}\) procede del término alemán \emph{Zahlen}, que significa ``números''. Los elementos de \(\mathbb{Z}\) se llaman números enteros.

Cada número natural positivo \(n\) tiene asociado un número entero negativo, denotado por \(-n\), que se llama su \emph{opuesto}. El número \(0\) es su propio opuesto, es decir,
\[
  -0=0.
\]

\begin{ejemplo}
Los números
\[
  -5,\quad -1,\quad 0,\quad 2,\quad 17
\]
son enteros. Los números
\[
  \frac{1}{2},\quad \sqrt{2},\quad \pi
\]
no son enteros.
\end{ejemplo}

\begin{convencion}[Inclusión de los naturales en los enteros]
Consideraremos siempre que
\[
  \mathbb{N}\subset \mathbb{Z}.
\]
Por tanto, todo número natural es también un número entero. Sin embargo, no todo número entero es natural: por ejemplo, \(-4\in\mathbb{Z}\), pero \(-4\notin\mathbb{N}\).
\end{convencion}

\section{Suma, producto y orden}

El conjunto \(\mathbb{Z}\) está dotado de dos operaciones fundamentales: la suma y el producto.

La suma de enteros se denota mediante el símbolo \(+\). Si \(a,b\in\mathbb{Z}\), entonces \(a+b\in\mathbb{Z}\). El producto de enteros se denota mediante yuxtaposición o mediante el símbolo \(\cdot\). Si \(a,b\in\mathbb{Z}\), entonces \(ab\in\mathbb{Z}\).

\begin{proposicion}[Propiedades básicas de la suma y del producto]
Para cualesquiera enteros \(a,b,c\in\mathbb{Z}\), se verifican las siguientes propiedades:
\begin{enumerate}
  \item \textbf{Asociatividad de la suma:}
  \[
    (a+b)+c=a+(b+c).
  \]
  \item \textbf{Conmutatividad de la suma:}
  \[
    a+b=b+a.
  \]
  \item \textbf{Elemento neutro de la suma:}
  \[
    a+0=a.
  \]
  \item \textbf{Elemento opuesto:}
  \[
    a+(-a)=0.
  \]
  \item \textbf{Asociatividad del producto:}
  \[
    (ab)c=a(bc).
  \]
  \item \textbf{Conmutatividad del producto:}
  \[
    ab=ba.
  \]
  \item \textbf{Elemento neutro del producto:}
  \[
    a\cdot 1=a.
  \]
  \item \textbf{Distributividad del producto respecto de la suma:}
  \[
    a(b+c)=ab+ac.
  \]
\end{enumerate}
\end{proposicion}

Estas propiedades indican que \(\mathbb{Z}\), con la suma y el producto, posee una estructura algebraica muy importante: es un anillo conmutativo con unidad. Esta terminología se estudiará con más detalle en cursos posteriores de álgebra.

\begin{advertencia}
Aunque todo entero tiene opuesto respecto de la suma, no todo entero distinto de cero tiene inverso respecto del producto dentro de \(\mathbb{Z}\). Por ejemplo, no existe ningún entero \(x\in\mathbb{Z}\) tal que
\[
  2x=1.
\]
Por tanto, \(2\) no tiene inverso multiplicativo en \(\mathbb{Z}\).
\end{advertencia}

El conjunto \(\mathbb{Z}\) también está ordenado. La relación de orden se denota por \(\leq\). Si \(a,b\in\mathbb{Z}\), escribimos
\[
  a\leq b
\]
para indicar que \(a\) es menor o igual que \(b\).

\begin{ejemplo}
Se tienen las desigualdades
\[
  -5<-2<0<3<8.
\]
En particular,
\[
  -5\leq -2,
  \qquad
  0\leq 3,
  \qquad
  -2\leq 8.
\]
\end{ejemplo}

El orden en los enteros es compatible con la suma: si \(a,b,c\in\mathbb{Z}\) y \(a\leq b\), entonces
\[
  a+c\leq b+c.
\]

También es compatible con el producto por enteros positivos: si \(a,b\in\mathbb{Z}\), \(c\in\mathbb{Z}\) y \(c>0\), entonces
\[
  a\leq b \quad \Longrightarrow \quad ac\leq bc.
\]
Si se multiplica por un entero negativo, el sentido de la desigualdad cambia.

\begin{ejemplo}
Como \(2<5\), al multiplicar por \(3>0\), obtenemos
\[
  6<15.
\]
En cambio, al multiplicar por \(-3<0\), obtenemos
\[
  -6>-15.
\]
\end{ejemplo}

\section{Valor absoluto}

El valor absoluto mide la distancia de un número entero al cero en la recta numérica.

\begin{definicion}[Valor absoluto]
Sea \(a\in\mathbb{Z}\). El valor absoluto de \(a\), denotado por \(|a|\), se define por
\[
  |a|=
  \begin{cases}
    a, & \text{si } a\geq 0,\\
    -a, & \text{si } a<0.
  \end{cases}
\]
\end{definicion}

De esta definición se deduce que \(|a|\) es siempre un número natural.

\begin{ejemplo}
Se tiene
\[
  |5|=5,
  \qquad
  |-5|=5,
  \qquad
  |0|=0.
\]
\end{ejemplo}

\begin{proposicion}[Propiedades elementales del valor absoluto]
Para cualesquiera enteros \(a,b\in\mathbb{Z}\), se verifican las siguientes propiedades:
\begin{enumerate}
  \item \(|a|\geq 0\).
  \item \(|a|=0\) si y solo si \(a=0\).
  \item \(|-a|=|a|\).
  \item \(|ab|=|a|\,|b|\).
  \item \(|a+b|\leq |a|+|b|\).
\end{enumerate}
\end{proposicion}

La última desigualdad se conoce como \emph{desigualdad triangular}. Su interpretación es que, si se avanza primero una distancia \(|a|\) y después una distancia \(|b|\), la distancia total al origen no puede ser mayor que la suma de ambas distancias.

\begin{ejemplo}
Si \(a=-3\) y \(b=5\), entonces
\[
  |a+b|=|-3+5|=|2|=2,
\]
mientras que
\[
  |a|+|b|=|-3|+|5|=3+5=8.
\]
Por tanto,
\[
  |a+b|\leq |a|+|b|.
\]
\end{ejemplo}

\section{División euclídea}

La división euclídea formaliza la división entera con resto. Es uno de los resultados fundamentales de la aritmética de los enteros y será esencial para estudiar la divisibilidad, el máximo común divisor y la aritmética modular.

\begin{teorema}[División euclídea]
Sean \(a,b\in\mathbb{Z}\), con \(b>0\). Entonces existen unos únicos enteros \(q,r\in\mathbb{Z}\) tales que
\[
  a=bq+r
\]
y
\[
  0\leq r<b.
\]
El número \(q\) se llama cociente y el número \(r\) se llama resto de la división de \(a\) entre \(b\).
\end{teorema}

En este enunciado, el entero \(a\) es el dividendo y el entero positivo \(b\) es el divisor.

\begin{ejemplo}
Si \(a=17\) y \(b=5\), entonces
\[
  17=5\cdot 3+2,
\]
con
\[
  0\leq 2<5.
\]
Por tanto, el cociente es \(q=3\) y el resto es \(r=2\).
\end{ejemplo}

\begin{ejemplo}
Si \(a=-17\) y \(b=5\), debemos encontrar enteros \(q\) y \(r\) tales que
\[
  -17=5q+r,
  \qquad
  0\leq r<5.
\]
Tomando \(q=-4\) y \(r=3\), se obtiene
\[
  -17=5(-4)+3=-20+3.
\]
Por tanto, el cociente es \(q=-4\) y el resto es \(r=3\).
\end{ejemplo}

\begin{advertencia}
Cuando el dividendo es negativo, el resto de la división euclídea sigue siendo no negativo. Por ejemplo, en la división de \(-17\) entre \(5\), el resto es \(3\), no \(-2\), porque el resto debe satisfacer
\[
  0\leq r<5.
\]
\end{advertencia}

\section*{Errores frecuentes}

\begin{erroresfrecuentes}
  \item Pensar que los números enteros negativos son números naturales. Con la convención de este manual, \(-1\notin\mathbb{N}\), aunque \(-1\in\mathbb{Z}\).
  \item Olvidar que al multiplicar una desigualdad por un número negativo se invierte el sentido de la desigualdad.
  \item Confundir el opuesto de un número, \(-a\), con su valor absoluto, \(|a|\).
  \item Creer que todos los enteros no nulos tienen inverso multiplicativo dentro de \(\mathbb{Z}\).
  \item Calcular incorrectamente el resto en una división euclídea con dividendo negativo.
\end{erroresfrecuentes}

\begin{seminario}
  \item Ejercicios de orden y valor absoluto.
  \item Aplicaciones de la división euclídea.
\end{seminario}

\begin{ejerciciospropuestos}
  \item Calcula el valor absoluto de los enteros \(-12\), \(0\), \(7\), \(-1\) y \(25\).
  \item Ordena de menor a mayor los enteros
  \[
    -3,
    \quad 5,
    \quad -8,
    \quad 0,
    \quad 2.
  \]
  \item Encuentra el cociente y el resto de la división euclídea de \(43\) entre \(6\).
  \item Encuentra el cociente y el resto de la división euclídea de \(-43\) entre \(6\).
  \item Demuestra que, para todo \(a\in\mathbb{Z}\), se tiene \(|-a|=|a|\).
  \item Da un ejemplo de dos enteros \(a,b\in\mathbb{Z}\) tales que
  \[
    |a+b|<|a|+|b|.
  \]
  \item Explica por qué no existe ningún entero \(x\in\mathbb{Z}\) tal que \(3x=1\).
\end{ejerciciospropuestos}

\resumen{Los números enteros amplían los números naturales al incorporar opuestos aditivos. El conjunto \(\mathbb{Z}\) permite resolver ecuaciones de la forma \(a+x=b\), está dotado de suma, producto y orden, y posee una herramienta aritmética fundamental: la división euclídea, que expresa todo entero como múltiplo de un divisor positivo más un resto acotado.}
